\section{The Puzzle}\label{sec:construction}

A \emph{hash chain} starts from a value and hashes it many times in a row. Each step needs the previous output, so even unlimited hardware needs the steps one after another. A single chain of $N$ steps would make a fine lock, except that its builder also needs $N$ sequential steps. The construction of Mahmoody, Moran and Vadhan~\cite{MMV11} fixes this: the sender builds $n$ short chains of $m$ steps at the same time and links them so that an opener can only run them in order. The opener's delay is $N = nm$ steps; the sender's wall time is $m$ steps on $n$ cores.

\subsection{Notation and model}\label{sec:prelims}\label{sec:notation}

The notation follows the specification (\texttt{SPEC.md} in the repository~\cite{leanSealRepo}). Integers are encoded big-endian, $a \cat b$ is concatenation, $\trunc{v}{b}$ is the first $b$ bytes of $v$, and $\oplus$ is bytewise XOR. A \emph{labelled hash} is $\THash(\mathit{label}, x) = \SHA(\mathsf{u8}(|\mathit{label}|) \cat \mathit{label} \cat x)$, and we write $\Hstart, \Hmask, \Hacc, \Hkey, \Hcommit$ for labelled hashes with the labels \hashlabel{start}, \hashlabel{mask} and so on. The \emph{chain step} for chain $i$ is
\[
  F_i(s) = \SHA\bigl(s \cat \salt \cat \mathsf{be32}(i)\bigr),
\]
where $\salt$ is 16 random bytes drawn for each puzzle. Its input is exactly 52 bytes, so one step is one SHA-256 compression, and no other hash in the format has an input of that length. A puzzle has $n$ chains of $m$ steps and $b$-byte links. The reference parameters are $n = 16$, $m = 2^{29}$ and $b = 24$: $N = 2^{33}$ steps, a 30-second delay at 3.5\,ns per step, and a 452-byte puzzle; we measured the one-minute puzzle with $m = 2^{30}$.

\label{sec:rom}We prove the delay in the random-oracle model~\cite{MMV11}: each chain step and each labelled hash is an independent random function that anyone can query, and an attacker is charged only for its queries. Time after publication is counted in \emph{oracle rounds}: all queries in a round are chosen before any answer from that round arrives. A bound in rounds therefore holds however much parallel hardware the attacker has; turning rounds into seconds takes the hardware assumption of \cref{sec:security}.

\subsection{Building and opening}\label{sec:algorithms}

\Cref{alg:build,alg:open} give the two algorithms, following Sections~4 and~5 of the specification with the byte encodings left out. The context $\ctx$ is a 32-byte encoding of the version, $b$, $n$, $m$ and the salt, and every hash other than the chain step includes it.

\begin{algorithm}[t]
  \caption{$\Build(n, m, b)$, run by the sender}\label{alg:build}
  \begin{algorithmic}[1]
    \State $\salt \gets$ 16 random bytes; \ $x_1, \dots, x_n \gets$ $b$ random bytes each
    \State $\ctx \gets (1, b, n, m, \salt)$
    \For{$i = 1, \dots, n$ \textbf{in parallel}}\label{line:build-for}
      \State $s_i \gets \Hstart(\ctx, i, x_i)$
      \State $y_i \gets F_i^m(s_i)$ \Comment{$m$ sequential steps}\label{line:build-chain}
    \EndFor
    \For{$i = 1, \dots, n - 1$}
      \State $z_i \gets \trunc{\Hmask(\ctx, y_i)}{b} \oplus x_{i+1}$ \Comment{hide the next link}\label{line:build-mask}
    \EndFor
    \State $a \gets \ctx$\label{line:build-acc}
    \State \algorithmicfor\ $i = 1, \dots, n$ \algorithmicdo\ $a \gets \Hacc(a, y_i)$
    \State $K \gets \Hkey(\ctx, a)$; \ $C \gets \Hcommit(\ctx, K)$\label{line:build-key}
    \State \Return puzzle $Z = (\ctx, C, x_1, z_1, \dots, z_{n-1})$ and key $K$
  \end{algorithmic}
\end{algorithm}

\begin{algorithm}[t]
  \caption{$\Open(Z)$, run by anyone}\label{alg:open}
  \begin{algorithmic}[1]
    \State parse $Z$ as $(\ctx, C, x_1, z_1, \dots, z_{n-1})$
    \State $x \gets x_1$; \ $a \gets \ctx$
    \For{$i = 1, \dots, n$} \Comment{in sequence}
      \State $y \gets F_i^m(\Hstart(\ctx, i, x))$ \Comment{$m$ sequential steps}
      \State $a \gets \Hacc(a, y)$ \Comment{lines~\ref{line:open-acc} and~\ref{line:open-mask} share one oracle round}\label{line:open-acc}
      \State \algorithmicif\ $i < n$ \algorithmicthen\ $x \gets \trunc{\Hmask(\ctx, y)}{b} \oplus z_i$ \Comment{unhide $x_{i+1}$}\label{line:open-mask}
    \EndFor
    \State $K \gets \Hkey(\ctx, a)$
    \State \algorithmicif\ $\Hcommit(\ctx, K) \neq C$ \algorithmicthen\ \Return $\bot$ \Comment{malformed puzzle}\label{line:open-check}
    \State \Return $K$
  \end{algorithmic}
\end{algorithm}

The sender knows every link, so it runs the $n$ chains of \cref{alg:build} in parallel. Each link after the first is published only masked, under a hash of the end of the chain before it. The key is a hash of all chain ends in order, and the commitment $C$ lets anyone check a claimed key with one hash. An opener without the key has only $x_1$: it runs chain~1, uses its end to unhide $x_2$, runs chain~2, and so on, $nm$ steps in sequence (\cref{fig:schedules}). The only way to start a chain early is to guess its link. The sender encrypts the transaction under $K$ with HMAC-SHA256 (\cref{sec:costs}) and keeps the puzzle and key until it needs them.

% Parallel building versus sequential opening (Figure 2 of the web version).
\begin{figure}[t]
  \centering
  \begin{tikzpicture}[
      x=1cm, y=1cm, font=\small,
      chain/.style={draw=accent, fill=accent!12, line width=0.5pt},
      note/.style={font=\footnotesize, text=black!60, align=left, anchor=west, inner xsep=0pt},
      head/.style={font=\small\bfseries, anchor=west},
    ]
    % Bar of m steps, with tick marks for the individual steps.
    \newcommand{\chainbar}[3]{% x0, y, width
      \draw[chain] (#1,#2) rectangle ++(#3,0.34);
      \foreach \t in {0.18,0.36,...,#3} {
        \draw[accent!45, line width=0.3pt] ({#1+\t},#2) -- ++(0,0.34);
      }
    }
    % ---------------- Sender ----------------
    \node[head] at (-0.3,1.25) {Sender: $n$ chains in parallel};
    \node[note] at (-0.3,0.85) {wall time $= m$ steps};
    \foreach \i/\y in {1/0, 2/-0.55, 3/-1.10, 4/-1.65} {
      \chainbar{0.35}{\y}{3.06}
      \node[anchor=east] at (0.3,\y+0.17) {$x_{\i}$};
      \node[anchor=west] at (3.45,\y+0.17) {$y_{\i}$};
    }
    \draw[decorate, decoration={brace, mirror, amplitude=4pt}]
      (0.35,-1.80) -- node[below=5pt] {$m$ hash steps} (3.41,-1.80);
    \node[note] at (-0.3,-2.75) {then publish $n - 1$ masked links,};
    \node[note] at (-0.3,-3.15) {$z_i = \trunc{\Hmask(\ctx, y_i)}{b} \oplus x_{i+1}$};
    % ---------------- Divider ----------------
    \draw[black!25] (4.45,1.45) -- (4.45,-3.75);
    % ---------------- Opener ----------------
    \node[head] at (4.75,1.25) {Opener: each chain needs the previous chain's end};
    \node[note] at (4.75,0.85) {wall time $= n \cdot m$ steps};
    \foreach \i/\x/\y in {1/5.20/0, 2/7.54/-0.55, 3/9.88/-1.10, 4/12.22/-1.65} {
      \chainbar{\x}{\y}{2.34}
    }
    \node[anchor=east] at (5.15,0.17) {$x_1$};
    \node[anchor=west] at (14.58,-1.48) {$y_4$};
    \foreach \i/\x/\y in {1/7.54/0, 2/9.88/-0.55, 3/12.22/-1.10} {
      \draw[-{Stealth[length=4pt]}, accent, line width=0.8pt] (\x,\y) -- (\x,\y-0.21);
      \node[anchor=west, text=accent, font=\footnotesize] at (\x+0.05,\y-0.08) {$\oplus\, z_{\i}$};
    }
    \draw[decorate, decoration={brace, mirror, amplitude=4pt}]
      (5.20,-1.80) -- node[below=5pt] {$n \cdot m$ hash steps; no parallelism helps} (14.56,-1.80);
    \node[note] at (4.75,-2.75) {Starting chain $i + 1$ early means guessing $x_{i+1}$,};
    \node[note] at (4.75,-3.15) {one of $2^{192}$ values for a 24-byte link.};
    \node[note] at (4.75,-3.55) {The key hashes every end $y_1, \dots, y_n$, so no chain can be left out.};
  \end{tikzpicture}
  \caption{The same $n$ chains, arranged two ways in time (drawn for $n = 4$). The sender knows every link and runs the chains side by side. An opener must run them one after another, because each chain's starting link is hidden under the end of the previous chain. The puzzle consists of the $n$ links, so its size grows with $n$ and not with the delay.}
  \label{fig:schedules}
\end{figure}


Four choices in the encoding matter for security. Every hash names its puzzle and chain, so work on another puzzle does not help. Masks are hashes of chain ends, not the ends themselves, so learning a link says nothing about the chain before it. The key depends on every chain's end, so a guessed link cannot be confirmed until the opener's own chains reach it. And the puzzle consists of $68 + nb$ bytes whatever the delay: the delay depends only on $N$.
